\chapter{The Composition Failure Theorem}
\label{ch:composition-failure}

\begin{quotation}
\noindent\textit{Synopsis.} This chapter presents the monograph's central
new result. The Composition Failure Theorem demonstrates that a sequence of
locally rational transformations can collectively produce a globally
insufficient outcome, with no stage behaving negligently, incompetently, or
irrationally: failure emerges because each stage optimizes for its own
local task while future tasks remain unknown. This formalizes a pattern
familiar from archival systems, institutions, scientific communication, and
information-processing pipelines, in which a distinction that seemed
irrelevant at one stage becomes essential later.
\end{quotation}

\newresult

\section{Making Local Admissibility Precise}

Before the theorem can be stated, the phrase it turns on --- a transition
being ``locally admissible with respect to its own stage's task''
--- needs a definition sharper than \cref{def:budget-discard}'s general
schema alone provides. \cref{def:budget-discard} says a distinction is
dropped when funding it exceeds budget or when it interferes too
strongly with a higher-priority distinction; it does not say how
priority itself is supposed to track a stage's task. The following
definition closes that gap, and everything downstream --- the theorem,
its corollaries, and the proof of Appendix~\ref{app:proof} --- depends on
exactly this closure.

\begin{definition}[Locally Admissible Transition]
\label{def:locally-admissible}
A transition $B_i$ is locally admissible with respect to $\tau_i$ if its
discard set $D_i$ is produced by a construction satisfying
\cref{def:budget-discard} in which every distinction required by $\tau_i$
(\cref{def:benign-loss}) is ranked strictly above every distinction not
required by $\tau_i$.
\end{definition}

This definition is deliberately weak in one respect and strong in
another. It is weak in that it does not choose between
\cref{con:greedy-discard} and \cref{con:optimal-discard}, or any other
construction satisfying \cref{def:budget-discard}: any of them counts as
locally admissible provided the priority ordering they operate on is
correctly aligned with $\tau_i$. It is strong in that this alignment is
not optional --- a transition that funds a distinction $\tau_i$ does not
need ahead of one it does need fails to be locally admissible, regardless
of how sophisticated the underlying discard construction is. A transition
satisfying \cref{def:locally-admissible} cannot be accused of
irrationality relative to $\tau_i$: it funds everything $\tau_i$ needs, in
full, before spending anything on distinctions $\tau_i$ does not need.
This is exactly the standard the theorem needs every stage to meet, since
the entire point of the result is that failure can occur even when no
stage falls short of it.

\section{Statement}

\begin{theorem}[Composition Failure]
\label{thm:composition-failure}
Let $A_0 \xrightarrow{B_0} \cdots \xrightarrow{B_{n-1}} A_n$ be a budget
relay (\cref{def:budget-relay}) in which each transition $B_i$ is locally
admissible (\cref{def:locally-admissible}) with respect to its own
stage's task $\tau_i$. Then there exist open task spaces $T$
(\cref{prop:task-unavailability}) containing some $\tau^* \notin \{\tau_0,
\ldots, \tau_{n-1}\}$ such that $A_n$ is insufficient
(\cref{def:sufficiency}) for $\tau^*$, even though no individual
transition was locally irrational relative to its own $\tau_i$.
\end{theorem}

Each clause of this statement is load-bearing, and it is worth checking
what happens if any one is dropped. Without \cref{def:budget-relay}'s
independence condition, the chain reduces to a single observer's economy
and \cref{thm:redistribution} governs it instead, as
Appendix~\ref{app:comparison} shows concretely. Without local
admissibility, the theorem would be uninteresting: of course an
incompetently-run relay can end up insufficient. Without $T$ being open,
$\tau^*$ could in principle have been anticipated and funded for, and the
failure would trace back to ordinary negligence rather than to the
structure of uncoordinated composition itself. The theorem's content is
precisely that none of these outs are available: the relay is
uncoordinated by definition, every stage did what it was supposed to do,
and the task that breaks it was not available to be planned for.

\subsection{Why No Stage Is Wrong}

It is worth stating this as plainly as possible, since it is the single
most commonly missed reading of the theorem. ``No stage is wrong'' does
not mean the outcome is somehow acceptable, or that no one should have
done anything differently in hindsight. It means that identifying which
stage to blame is a type error: blame presupposes a standard the blamed
party fell short of, and every stage in \cref{thm:composition-failure}'s
hypothesis met its own standard exactly. Stage $i$'s standard is
$\tau_i$; stage $i$'s transition $B_i$ is locally admissible with
respect to $\tau_i$ by hypothesis; there is no further standard stage
$i$ could have been held to without already knowing $\tau^*$, which
\cref{prop:task-unavailability} rules out. Asking ``which stage failed''
is asking a question with a false presupposition, in the same way that
asking which digit of a correctly computed sum is wrong is a false
presupposition when the sum is correct. The theorem's discomfort, for
readers who feel there ought to be someone to hold responsible, is not a
flaw in the theorem; it is the theorem's content, restated as a feeling.

\subsection{Relation to Arrow-Type Impossibility Results}

Readers familiar with social choice theory may notice a family
resemblance to Arrow's impossibility theorem, which shows that no voting
rule aggregating individual preferences can simultaneously satisfy a
short list of independently reasonable fairness conditions. The
resemblance is worth naming and then bounding carefully, since the two
results are not the same theorem in different clothing. Both share the
structure of showing that a collection of locally reasonable
requirements --- individually rational voters, individually rational
relay stages --- cannot be jointly satisfied by any aggregation
procedure meeting a further, seemingly modest condition (independence of
irrelevant alternatives, in Arrow's case; uncoordinated composition, in
this monograph's). Where they differ is in what is being aggregated and
against what standard: Arrow's theorem is about preference aggregation
under a fixed, closed set of alternatives and voters; \cref{thm:composition-failure}
is about distinction preservation across an open task space, and its
failure mode is diachronic in a way Arrow's is not, since it depends on
a later task appearing that no stage could have voted on in the first
place because it did not yet exist. The comparison is offered as an
aid to intuition for readers who already know Arrow's theorem, not as a
claim that one result specializes to or generalizes the other; no such
reduction is attempted here.

\section{Proof Strategy}

The full proof is given in Appendix~\ref{app:proof}; this section
previews its shape so the construction does not appear unmotivated when
read there.

\begin{proofsketch}
The proof constructs an explicit witness relay using the five-distinction
instance introduced in \cref{sec:worked-relay} and reused in
\cref{sec:redistribution-worked}, together with a family of
\emph{recovery tasks} $\tau_d$ --- one for each named distinction $d$,
asking only whether $d$ holds in the source
(\cref{def:recovery-task}) --- which serves as the open task space $T$.
A stage whose required distinctions exactly, or near-exactly, exhaust its
local budget is shown (\cref{lem:exhaustion}) to drop every non-required
distinction \emph{regardless of which construction satisfying
\cref{def:budget-discard} is used}, which is what allows the witness to
be locally admissible in the strong sense of \cref{def:locally-admissible}
without leaving room for a more clever discard policy to have rescued the
dropped distinction. Some distinction $d^*$ -- concretely, $d_5$ in the
worked witness -- is required by neither stage's task and is dropped at
the earliest stage as a direct consequence. By
\cref{prop:monotonicity}, $d^*$ remains unavailable at every later stage
absent root recovery (\cref{prop:root-recovery}), which the witness does
not invoke. Taking $\tau^* = \tau_{d^*}$ then gives insufficiency
directly, by \cref{lem:recovery-sufficiency}: a rendering is sufficient
for a recovery task exactly when it has not discarded the distinction in
question, and $A_n$ has.
\end{proofsketch}

This preview differs in one respect from an earlier, less precise sketch
of the same argument, which described $d^*$ as dropped because it
``interferes with a distinction $\tau_i$ does require.'' The actual
mechanism proved in Appendix~\ref{app:proof} is sharper and does not rely
on interference at all: $d^*$ is dropped because, after funding every
required distinction in full, no budget capable of funding $d^*$ remains
--- a conclusion that holds regardless of interference structure and
regardless of which construction satisfying \cref{def:budget-discard}
produced the allocation. The interference-based mechanism remains
available in principle for other witnesses, but it is not what the
canonical witness of this monograph uses, and the sketch above is
corrected accordingly.

\section{Corollaries}

\begin{corollary}[Local Improvement Does Not Guarantee Global Sufficiency]
\label{cor:local-improvement}
Improving any single stage's local budget cannot, by itself, eliminate
all instances of the failure mode exhibited in
\cref{thm:composition-failure}.
\end{corollary}

\begin{proofsketch}
Verified directly against the witness in Appendix~\ref{app:proof}:
increasing $\beta_1$ cannot restore a distinction already absent from
$A_1$ by the time stage $1$ acts, since \cref{prop:monotonicity} forbids
recovery absent root access. Only an improvement at the specific stage
where the loss occurred could have prevented it, and even that requires
knowing, at that stage, that the eventually-relevant task would need the
distinction in question -- precisely the knowledge
\cref{prop:task-unavailability} shows is generally unavailable.
\end{proofsketch}

\begin{corollary}[No Free Lunch for Relays]
\label{cor:no-free-lunch}
A budget relay with $n$ uncoordinated stages has cumulative discard
(\cref{def:cumulative-discard}) that is non-decreasing in $n$
(\cref{prop:monotonicity}), realized with equality across arbitrarily
many additional pass-through stages
(\cref{prop:extension}); the coordination mechanisms of
Chapter~\ref{ch:coordination} can narrow, but not eliminate, the
resulting insufficiency (\cref{prop:no-coordination-universal}).
\end{corollary}

\begin{proofsketch}
The non-decrease is \cref{prop:monotonicity} directly, and
Appendix~\ref{app:proof} shows this bound is achieved with equality
rather than merely respected, since pass-through stages contribute no
further discard. This corollary should not be read as a claim that
cumulative discard \emph{risk} grows in some stronger, quantitative
sense across relays in general; that stronger claim is not established
here and is recorded as open in Appendix~\ref{app:open-questions}. What
is established is the existence half only: discard does not shrink, and
can accumulate to the full extent this chapter's witness demonstrates.
\end{proofsketch}

\section{What the Theorem Does Not Say}

It is worth listing plainly what \cref{thm:composition-failure} does not
claim, since a result this general invites over-reading in several
directions at once. It does not say that every budget relay fails; the
homogeneous-task chains of \cref{prop:homogeneous-chain}
(Chapter~\ref{ch:proofs}) are relays that provably do not, and closed
hidden-state domains under augmentation (\cref{prop:closed-sufficiency})
are a further case where sufficiency is achieved outright. It does not
say that coordination is worthless; Chapter~\ref{ch:coordination}'s
mechanisms measurably narrow the failure mode even though
\cref{prop:no-coordination-universal} shows none of them close it
entirely. It does not identify a rate at which insufficiency grows,
only that it does not shrink (\cref{cor:no-free-lunch}), and the
sharper, quantitative version of that claim is explicitly left open
rather than implied. It does not say anything at all about whether a
given real institution, pipeline, or memory system is currently
suffering from this failure mode; the theorem establishes that the
failure mode is possible under stated hypotheses, not that it is
occurring in any particular case, which would require checking those
hypotheses against the case in question. And it does not require malice,
incompetence, or even error at any stage, which is the point the
preceding two sections existed to make and which this list exists to
keep from being forgotten once the theorem is put to use.

\section{Relation to the Redistribution Theorem}

\cref{thm:composition-failure} is not a special case of
\cref{thm:redistribution}, and is not derivable from it without the
budget-relay machinery of Chapter~\ref{ch:budget-relay}. The
Redistribution Theorem concerns simultaneous competition for one
observer's budget; the Composition Failure Theorem concerns a single
distinction's fate across a sequence of independent, non-communicating
budgets. Appendix~\ref{app:comparison} develops this contrast concretely
rather than only structurally: pooling the witness's two stage budgets
into a single shared total is shown there
(\cref{prop:pooling-restores}) not merely to make
\cref{thm:redistribution} applicable again, but to actually succeed at
funding the distinction the separated relay was forced to drop, using
the identical total nominal resource. The insufficiency this chapter's
theorem exhibits is therefore a consequence specifically of
\cref{def:budget-relay}'s prohibition on transferring unspent capacity
across stages, not of any absolute scarcity of resource -- confirming
that the two theorems occupy genuinely disjoint settings rather than one
subsuming the other.
